\newcommand{\TabOptimisationParameters}{%
\begin{table}[tb]%
\begin{tabular}{ll}%
	\hline
	Region for optimisation & Approx.\ No.\  of parameters \\
	\hline
	Production target dimensions and location & $3+3$ \\
	Torus1 dipole field strength & 1 \\
	Torus2 dipole field strength & 1 \\
	Muon beam collimator shapes, position, and material & $3+1+1$ \\
	Stopping target shape and location & $4+3$ \\
	Beam blocker position, form, and material& $3+3+1$ \\
	Electron spectrometer dipole field strength& 1 \\
	DIO blockers in the spectrometer & $4$ \\
	\hline
	Approx. total number of parameters& 32 \\
	\hline
\end{tabular}%
	\caption{\tablabel{optimisation:possible-parameters}%
	Aspects of the experiment that can be optimised and estimates for the number of parameters that define each aspect.
	In the case of the target, beam blocker, and collimator shapes the number of parameters is only approximate; crudely speaking there is at least a width, length and height but in principle one could have a very irregular shape that cannot be parametrised by only three numbers, for example, shapes that change as a function of distance along the beamline.
}%
\end{table}%
}

\newcommand{\FigOptimProdTgtLength}{
\begin{figure}[pt]
\centering
\subfloat[][\figlabel{optimisation:ProdTgtSec:Length:Momentum:Muons}Muons]{
\includegraphics[width=0.48\textwidth,trim=0 0 0 1.5cm,clip]{figs/optimisation/ProdTgtGeom/Length_mu-minus_momentum}}
\subfloat[][\figlabel{optimisation:ProdTgtSec:Length:Momentum:Pions}Pions]{
\includegraphics[width=0.48\textwidth,trim=0 0 0 1.5cm,clip]{figs/optimisation/ProdTgtGeom/Length_pi-minus_momentum}}
\caption{\figlabel{optimisation:ProdTgtSec:Length:Momentum}
Change to momentum distributions at the entrance to the first 90 degrees of the bent muon beam solenoid for different target lengths.
Target length is given as half-length which is the Geant4 convention.  
}
\end{figure}
\begin{figure}[pt]
\centering
\subfloat[][\figlabel{optimisation:ProdTgtSec:Length:Integral:Muons}Muons]{
\includegraphics[width=0.48\textwidth,trim=0 0 2cm 2.8cm,clip]{figs/optimisation/ProdTgtGeom/Length_mu-minus_integral_toZero}}
\subfloat[][\figlabel{optimisation:ProdTgtSec:Length:Integral:Pions}Pions]{
\includegraphics[width=0.48\textwidth,trim=0 0 2cm 2.8cm,clip]{figs/optimisation/ProdTgtGeom/Length_pi-minus_integral_toZero}}
\caption{\figlabel{optimisation:ProdTgtSec:Length:Integral}
Integrated muon and pion yields up to a certain momentum at the entrance to the first 90 degrees of the bent muon beam solenoid as a function of target length.
}
\end{figure}
\begin{figure}[pt]
\centering
\subfloat[][\figlabel{optimisation:ProdTgtSec:Length:IntegralRatio:Muons}Muons]{\includegraphics[width=0.48\textwidth,trim=0 0 2cm 2.8cm,clip]{figs/optimisation/ProdTgtGeom/Length_mu-minus_integral_ratios}}
\subfloat[][\figlabel{optimisation:ProdTgtSec:Length:IntegralRatio:Pions}Pions]{\includegraphics[width=0.48\textwidth,trim=0 0 2cm 2.8cm,clip]{figs/optimisation/ProdTgtGeom/Length_pi-minus_integral_ratios}}
\caption{\figlabel{optimisation:ProdTgtSec:Length:IntegralRatio}
Change in the momentum distribution of muons and pions at the entrance to the first 90 degrees of the bent muon beam solenoid as a function of target length.
}
\end{figure}
}

\newcommand{\FigOptimProdTgtRad}{
\begin{figure}[pt]
\centering
\subfloat[][\figlabel{optimisation:ProdTgtSec:Radius:Momentum:Muons}Muons]{\includegraphics[width=0.48\textwidth,trim=0 0 0 1.5cm,clip]{figs/optimisation/ProdTgtGeom/Radius_mu-minus_momentum}}
\subfloat[][\figlabel{optimisation:ProdTgtSec:Radius:Momentum:Pions}Pions]{\includegraphics[width=0.48\textwidth,trim=0 0 0 1.5cm,clip]{figs/optimisation/ProdTgtGeom/Radius_pi-minus_momentum}}
\caption{
\figlabel{optimisation:ProdTgtSec:Radius:Momentum}
Change to momentum distributions at the entrance to the first 90 degrees of the bent muon beam solenoid for different target radii.
}
\end{figure}
\begin{figure}[pt]
\centering
\subfloat[][\figlabel{optimisation:ProdTgtSec:Radius:Integral:Muons}Muons]{\includegraphics[width=0.48\textwidth,trim=0 0 2cm 2.8cm,clip]{figs/optimisation/ProdTgtGeom/Radius_mu-minus_integral_toZero}}
\subfloat[][\figlabel{optimisation:ProdTgtSec:Radius:Integral:Pions}Pions]{\includegraphics[width=0.48\textwidth,trim=0 0 2cm 2.8cm,clip]{figs/optimisation/ProdTgtGeom/Radius_pi-minus_integral_toZero}}
\caption{\figlabel{optimisation:ProdTgtSec:Radius:Integral}
Integrated muon and pion yields up to a certain momentum at the entrance to the first 90 degrees of the bent muon beam solenoid as a function of target radius.
}
\end{figure}
\begin{figure}[pt]
\centering
\subfloat[][\figlabel{optimisation:ProdTgtSec:Radius:IntegralRatio:Muons}Muons]{\includegraphics[width=0.48\textwidth,trim=0 0 2cm 2.8cm,clip]{figs/optimisation/ProdTgtGeom/Radius_mu-minus_integral_ratios}}
\subfloat[][\figlabel{optimisation:ProdTgtSec:Radius:IntegralRatio:Pions}Pions]{\includegraphics[width=0.48\textwidth,trim=0 0 2cm 2.8cm,clip]{figs/optimisation/ProdTgtGeom/Radius_pi-minus_integral_ratios}}
\caption{\figlabel{optimisation:ProdTgtSec:Radius:IntegralRatio}
Change in the momentum distribution of muons and pions at the entrance to the first 90 degrees of the bent muon beam solenoid as a function of target radius.
}
\end{figure}
}

\newcommand{\FigOptimProdTgtFinal}{
\begin{figure}[pt]
\centering
\subfloat[][\figlabel{optimisation:ProdTgtSec:Final:Integral:Muons}Muons]{\includegraphics[width=0.48\textwidth,trim=0 0 0 1.5cm,clip]{figs/optimisation/ProdTgtGeom/OptimalLengthRadius_mu-minus_integral_toZero.pdf}}
\subfloat[][\figlabel{optimisation:ProdTgtSec:Final:Integral:Pions}Pions]{\includegraphics[width=0.48\textwidth,trim=0 0 0 1.5cm,clip]{figs/optimisation/ProdTgtGeom/OptimalLengthRadius_pi-minus_integral_toZero.pdf}}
\caption{\figlabel{optimisation:ProdTgtSec:Final:Integral}
Variation in muon and pion yields as a function of target radius when the total target length is set to the optimised value of 32~cm.
Despite the longer target length the optimal radius is still 1~cm.
}
\end{figure}
}

\newcommand{\FigOptimProdTgtComparePhases}{
\begin{figure}[pt]
\centering
\includegraphics[width=0.9\textwidth,trim=0 0 0 1.5cm,clip]{figs/optimisation/ProdTgtGeom/Plot_compare_phase_1and2.pdf}
\caption{\figlabel{optimisation:ProdTgtSec:Phase1vs2}
Comparison of the muon and pion yields per POT for \phaseI and \phaseII.
The difference arises from the change of target material between the phases.
}
\end{figure}
}
\newcommand{\FigOptimMuBeamDipoleMuStops}{
\begin{figure}[bt]
\centering
%	\fbox{
\includegraphics[width=0.85\textwidth,trim=0 0.5cm 0 1.0cm,clip]{figs/optimisation/MuonBeamDipoles/Tidied_stopped_muons.pdf}
%}
\caption{\figlabel{optim:muBeamDipole:stoppedMu}
	Muon stopping rate as a function of the two dipole field strengths (given relative to the \phaseI design specification).
	A clear anti-correlation is visible which is discussed in the text.
}
\end{figure}
}

\newcommand{\FigOptimMuBeamDipolePiStops}{
\begin{figure}[t]
\centering
%	\fbox{
\includegraphics[width=0.85\textwidth,trim=0 0.5cm 0 1.0cm,clip]{figs/optimisation/MuonBeamDipoles/Tidied_stopped_pions.pdf}
%}
\caption{\figlabel{optim:muBeamDipole:stoppedPi}
	Pion stopping rate as a function of the two dipole field strengths (given relative to the \phaseI design specification).
At the level of statistics used to generate each point, no clear trend is obvious.
Empty squares are those where no pions stopped in the run.
}
\end{figure}
}

\newcommand{\FigOptimMuBeamDipoleMuDispersion}{
\begin{figure}[t]
\centering 
\subfloat[][\figlabel{optim:MuBeamDipole:MuDispersion:Entry}Torus1 Entry]{\includegraphics[width=0.45\textwidth,trim=0 0.9cm 0 1.9cm,clip]{figs/optimisation/MuonBeamCollimators/Tidied_Muon_dispersion_at_entrance.pdf}}
\subfloat[][\figlabel{optim:MuBeamDipole:MuDispersion:Exit}Torus2 Exit]  {\includegraphics[width=0.45\textwidth,trim=0 0.9cm 0 1.9cm,clip]{figs/optimisation/MuonBeamCollimators/Tidied_Muon_dispersion_at_exit.pdf}}
\caption{\figlabel{optim:MuBeamDipole:MuDispersion}
Dispersive effect of the 180\degree bent transport solenoid and dipole field on muons.
No collimating material is yet included, so the high-energy muons being removed is due purely to the beam-pipe itself.
}
\end{figure}
}

\newcommand{\FigOptimMuBeamCollimMuonPaths}{
\begin{figure}[p]
\centering 
	\subfloat[][\figlabel{optim:MuBeamCollim:Beamline:All}All Muons]          {\includegraphics[width=1\textwidth,trim=18cm 1.0cm 26cm 1cm,clip]{figs/optimisation/MuonBeamCollimators/Tidied_All_muons_wGeom.png}}\\
\subfloat[][\figlabel{optim:MuBeamCollim:Beamline:Stopped}Stopped Muons]          {\includegraphics[width=1\textwidth,trim=18cm 1.0cm 26cm 1cm,clip]{figs/optimisation/MuonBeamCollimators/Tidied_Stopped_muons_wGeom.png}}\\
\subfloat[][\figlabel{optim:MuBeamCollim:Beamline:HighP}Muons with $p>70$ MeV/c around the stopping target]          {\includegraphics[width=1\textwidth,trim=18cm 1.0cm 26cm 1cm,clip]{figs/optimisation/MuonBeamCollimators/Tidied_HighP_muons_wGeom.png}}\\
\subfloat[][\figlabel{optim:MuBeamCollim:Beamline:Diff}Stopped $\mu-$High-$p$ $\mu$]{\includegraphics[width=1\textwidth,trim=18cm 1.0cm 26cm 1cm,clip]{figs/optimisation/MuonBeamCollimators/Tidied_Where_to_collimate_wGeom.png}}
\caption{\figlabel{optim:MuBeamCollim:Beamline}
The heights of muons as they pass along the beamline.  
	\protect\subref{fig:optim:MuBeamCollim:Beamline:All} The path of all muons.
	\protect\subref{fig:optim:MuBeamCollim:Beamline:Stopped}: The paths of muons that stop in the target.
	\protect\subref{fig:optim:MuBeamCollim:Beamline:HighP}: The heights of muons with momentum greater than 70 MeV/c when they enter the region around the stopping target.  These could potentially decay in flight to give electrons with 100 MeV/c or greater.
	\protect\subref{fig:optim:MuBeamCollim:Beamline:Diff}: The difference between plot \protect\subref{fig:optim:MuBeamCollim:Beamline:Stopped} and plot \protect\subref{fig:optim:MuBeamCollim:Beamline:HighP}.
	Regions in dark blue would give the greatest impact in removing high-momentum muons whilst leave the stopping muons untouched.
	These plots should be compared to those of \fig{optim:MuBeamCollim:BeamWColl} once collimators have been introduced.
}
\end{figure}
}

\newcommand{\FigOptimMuBeamCollimTransverseSep}{
\begin{figure}[t]
\centering 
\subfloat[][\figlabel{optim:MuBeamCollim:TransverseSep:TS1Entry}At 0\degree]  {\includegraphics[height=0.35\textheight,trim=0.0cm 0.8cm 1.3cm 1.9cm,clip]{figs/optimisation/MuonBeamCollimators/MuonTransversePos_Torus1Tor1FirstPoint}}
\subfloat[][\figlabel{optim:MuBeamCollim:TransverseSep:TS3}At 90\degree] {\includegraphics[height=0.35\textheight,trim=1.7cm 0.8cm 1.3cm 1.9cm,clip]{figs/optimisation/MuonBeamCollimators/MuonTransversePos_Torus2Monitor_0}}
\subfloat[][\figlabel{optim:MuBeamCollim:TransverseSep:TS2Exit}At 180\degree]{\includegraphics[height=0.35\textheight,trim=1.7cm 0.8cm 1.3cm 1.9cm,clip]{figs/optimisation/MuonBeamCollimators/MuonTransversePos_StopTgtSecMonitor_0}}
\caption{\figlabel{optim:MuBeamCollim:TransverseSep}
The separation between stopping and dangerous muons.
The separation is largest at the exit (180\degree), reasonable at the entrance (0\degree), and smallest around the mid-point (90\degree).
}
\end{figure}
}

\newcommand{\FigOptimMuBeamCollimMuonPathsWColl}{
\begin{figure}[ph]
\centering 
\subfloat[][\figlabel{optim:MuBeamCollim:BeamWColl:All}All Muons]                  {\includegraphics[width=1\textwidth,trim=18cm 1.0cm 26cm 1cm,clip]{figs/optimisation/MuonBeamCollimators/Tidied_WColl_AllMuons_WGeom.png}}\\
\subfloat[][\figlabel{optim:MuBeamCollim:BeamWColl:Stopped}Stopped Muons]          {\includegraphics[width=1\textwidth,trim=18cm 1.0cm 26cm 1cm,clip]{figs/optimisation/MuonBeamCollimators/Tidied_WColl_StoppedMuons_WGeom.png}}\\
\subfloat[][\figlabel{optim:MuBeamCollim:BeamWColl:HighP}Muons with $p>70$ MeV/c around the stopping target]          {\includegraphics[width=1\textwidth,trim=18cm 1.0cm 26cm 1cm,clip]{figs/optimisation/MuonBeamCollimators/Tidied_WColl_HighPMuons_WGeom.png}}\\
\caption{\figlabel{optim:MuBeamCollim:BeamWColl}
The heights of muons as they pass along the beamline.  
	\protect\subref{fig:optim:MuBeamCollim:Beamline:All} The path of all muons.
	\protect\subref{fig:optim:MuBeamCollim:Beamline:Stopped}: The paths of muons that stop in the target.
	\protect\subref{fig:optim:MuBeamCollim:Beamline:HighP}: The heights of muons with momentum greater than 70 MeV/c when they enter the region around the stopping target.  These could potentially decay in flight to give electrons with 100 MeV/c or greater.
	These plots should be compared to those of \fig{optim:MuBeamCollim:Beamline} before collimators were introduced, where it is clear how well the dangerous muons are being suppressed.
}
\end{figure}
}

\newcommand{\FigOptimMuBeamCollimTorusOne}{
\begin{figure}[t]
\centering 
\subfloat[][\figlabel{optim:MuBeamCollim:Torus1:perPOT}Particles Surviving per POT]{\includegraphics[width=0.5\textwidth,trim=0.8cm 0.8cm 0.6cm 1.9cm,clip]{figs/optimisation/MuonBeamCollimators/Survived_Coll1_unNormalised-log.pdf}}
\subfloat[][\figlabel{optim:MuBeamCollim:Torus1:fraction}Fraction Surviving Collimator]{\includegraphics[width=0.5\textwidth,trim=0.8cm 0.8cm 0.6cm 1.9cm,clip]{figs/optimisation/MuonBeamCollimators/Survived_Coll1_Normalised-lin.pdf}}
\caption{\figlabel{optim:MuBeamCollim:Torus1}
The effect of changing the height of the collimator in Torus1 on the particle distributions.
}
\end{figure}
}

\newcommand{\FigOptimMuBeamCollimTorusTwoPerPOT}{
\begin{figure}[t]
\centering 
\includegraphics[width=0.95\textwidth,trim=0.3cm 0.3cm 0.8cm 0.1cm,clip]{figs/optimisation/MuonBeamCollimators/Survived_Coll2_unNormalised-lin.pdf}
\caption{\figlabel{optim:MuBeamCollim:Torus2:perPOT}
The number of particles reaching the end of the Torus2 solenoid per POT for different heights of both collimators in Torus1 and Torus2.
}
\end{figure}
}


\newcommand{\FigOptimMuBeamCollimTorusTwoFraction}{
\begin{figure}[t]
\centering 
\includegraphics[width=0.8\textwidth,trim=0.3cm 0.3cm 0.8cm 0.1cm,clip]{figs/optimisation/MuonBeamCollimators/Survived_Coll2_unNormalised-lin.pdf}
\caption{\figlabel{optim:MuBeamCollim:Torus2:fraction}
	The number of particles reaching the end of the Torus2 solenoid relative to the number that enter the Torus1 solenoid (\ie the survival probability) for different heights of both collimators in Torus1 and Torus2.
}
\end{figure}
}

\newcommand{\FigOptimMuBeamCollimTorusTwoContours}{
\begin{figure}[bt]
\centering 
%\fbox{
\includegraphics[width=0.75\textwidth,trim=0.1cm 0.3cm 2.8cm 1.1cm,clip]{figs/optimisation/MuonBeamCollimators/Survived_Coll2_StoppedVsHighP-Muons.pdf}
%}
\caption{\figlabel{optim:MuBeamCollim:Torus2:contours}
Contours showing 2.5 percentage point changes to the stopping (blue) and dangerous (red) muon flux, as a function of the collimator heights.
100\% acceptance is found in the bottom right corner.
%For example, for collimator heights within the first blue contour towards the bottom-right corner, less than 2.5\% of stopped muons are lost.
}
\end{figure}
}

\newcommand{\FigOptimESTDipoleBeamHeightTwoD}{
\begin{figure}[tp]
\centering 
	\subfloat[][\figlabel{optim:ESTDipole:Beam:0.0}No Dipole]  {\includegraphics[width=1\textwidth,trim=4cm 0.5cm 9.5cm 0.7cm,clip]{figs/optimisation/EST_dipole/Tidied_signal_height-dipole_00}}\\
\subfloat[][\figlabel{optim:ESTDipole:Beam:0.1}0.1~T Dipole]{\includegraphics[width=1\textwidth,trim=4cm 0.5cm 9.5cm 0.7cm,clip]{figs/optimisation/EST_dipole/Tidied_signal_height-dipole_10}}\\
\subfloat[][\figlabel{optim:ESTDipole:Beam:0.2}0.2~T Dipole]{\includegraphics[width=1\textwidth,trim=4cm 0.5cm 9.5cm 0.7cm,clip]{figs/optimisation/EST_dipole/Tidied_signal_height-dipole_20}}\\
\caption{\figlabel{optim:ESTDipole:Beam}
The heights of signal electrons for different dipole field values.
}
\end{figure}
}

\newcommand{\FigOptimESTDipoleBeamHeightMean}{
\begin{figure}[tb]
\centering 
%	\fbox{
\includegraphics[width=0.8\textwidth,trim=1.0cm 1.3cm 3.8cm 0.9cm,clip]{figs/optimisation/EST_dipole/Tidied_NoShift-Height}
%}
\caption{\figlabel{optim:ESTDipole:MeanHeight}
Mean height of signal electrons for different values of the dipole field strength.
}
\end{figure}
}

\newcommand{\FigOptimESTDipoleBeamFluxMean}{
\begin{figure}[tb]
\centering 
\includegraphics[width=0.8\textwidth,trim=1.0cm 0.3cm 3.8cm 1.8cm,clip]{figs/optimisation/EST_dipole/Tidied_NoShift-Flux}
\caption{\figlabel{optim:ESTDipole:MeanFlux}
Survival probability for signal electrons as a function of the distance along the beamline for different values of the electron spectrometer's dipole field strengths.
}
\end{figure}
}

\newcommand{\FigOptimESTDipoleAcceptanceVsDipole}{
\begin{figure}[tb]
\centering 
\includegraphics[width=0.6\textwidth,trim=0.3cm 0.3cm 1.9cm 1.1cm,clip]{figs/optimisation/EST_dipole/Tidied_acceptance}
\caption{\figlabel{optim:ESTDipole:acceptance}
Geometric acceptance into the StrECAL detector as a function of the dipole field strength over the electron spectrometer.
}
\end{figure}
}

\newcommand{\FigOptimStopTgtPosMuStops}{
\begin{figure}[bt]
\centering 
%	\fbox{
\includegraphics[width=0.7\textwidth,trim=2.0cm 0.05cm 0.9cm 0.3cm,clip]{figs/optimisation/StopTgtPosition/Tidied_MuonStoppingRate.pdf}
%}
\caption{\figlabel{optim:StopTgtPos:MuStops}
Muon stopping rate per \ac{POT} for different target positions.
The linear behaviour arises from the reduced field strength and fixed target radius such that fewer muons impact the target as it is moved downstream.
}
\end{figure}
}

\newcommand{\FigOptimStopTgtPosSensitivitySpect}{
\begin{figure}[tb]
\centering 
%	\fbox{
\includegraphics[width=0.9\textwidth,,trim=1.3cm 0.2cm 0.7cm 0.6cm,clip]{figs/optimisation/StopTgtPosition/Tidied_-AcceptedMomentum.pdf}
%}
\caption{\figlabel{optim:StopTgtPos:AcceptedMomSpect}
The momentum dependence of the electron acceptance into the detector for different target positions.
The spectrum for each target position is normalised to the muon stopping rate for that position, such that each curve shows the sensitivity to electrons of that momentum.
}
\end{figure}
}

\newcommand{\FigOptimStopTgtPosSensitivityNoBeamBlock}{
\begin{figure}[tb]
\centering 
%	\fbox{
\includegraphics[width=0.9\textwidth,trim=1.3cm 0.2cm 0.7cm 0.6cm,clip]{figs/optimisation/StopTgtPosition/Tidied_NoBeamBlocker-AcceptedMomentum.pdf}
%}
\caption{\figlabel{optim:StopTgtPos:AcceptedMomSpectNoBeamBlock}
The momentum dependence of the electron acceptance into the detector for different target positions when the beam blocker is removed.
}
\end{figure}
}

\newcommand{\FigOptimStopTgtPosSensitivityIntegral}{
\begin{figure}[tb]
\centering 
%	\fbox{
\subfloat[][\figlabel{optim:StopTgtPos:AcceptIntegral:Adjacent}Sensitivity]{%
\includegraphics[width=0.49\textwidth,trim=0.5cm 0.5cm 0.3cm 1.9cm,clip]{figs/optimisation/StopTgtPosition/Tidied_-AcceptedMomentum-Integrated_adjacent.pdf}}
\subfloat[][\figlabel{optim:StopTgtPos:AcceptIntegral:Ratio}Change in Shape]{%
\includegraphics[width=0.49\textwidth,trim=0.5cm 0.5cm 0.3cm 1.9cm,clip]{figs/optimisation/StopTgtPosition/Tidied_-AcceptedMomentum-Integrated_ratios.pdf}}
%}
\caption{\figlabel{optim:StopTgtPos:AcceptIntegral}
\protect\subref{fig:optim:StopTgtPos:AcceptIntegral:Adjacent} 
The variation in sensitivity (acceptance $\times$ stopping rate) to electrons with different momenta as a function of the target position with respect to the nominal location. 
The darkest red line towards the top of the plot represents the sensitivity to signal, and it is that line that should therefore be maximised.
\protect\subref{fig:optim:StopTgtPos:AcceptIntegral:Ratio} 
The change in the shape of the acceptance vs. momentum spectrum as a function of the stopping target location.
}
\end{figure}
}

\newcommand{\FigOptimStopTgtPosHeights}{
\begin{figure}[p]
\centering 
\subfloat[][\figlabel{optim:StopTgtPos:Height:42.5}Electrons from 40 to 45 MeV/c]{%
\includegraphics[width=0.95\textwidth,trim=1.15cm 0.05cm 0.3cm 0.180cm,clip]{figs/optimisation/StopTgtPosition/WithBeamBlocker-Height-VaryShifts-Momentum_42-5.pdf}}%
\\\subfloat[][\figlabel{optim:StopTgtPos:Height:62.5}Electrons from 60 to 65 MeV/c]{%
\includegraphics[width=0.95\textwidth,trim=1.15cm 0.05cm 0.3cm 0.180cm,clip]{figs/optimisation/StopTgtPosition/WithBeamBlocker-Height-VaryShifts-Momentum_62-5.pdf}}%
\\\subfloat[][\figlabel{optim:StopTgtPos:Height:82.5}Electrons from 80 to 85 MeV/c]{%
\includegraphics[width=0.95\textwidth,trim=1.15cm 0.05cm 0.3cm 0.180cm,clip]{figs/optimisation/StopTgtPosition/WithBeamBlocker-Height-VaryShifts-Momentum_82-5.pdf}}%
\\\subfloat[][\figlabel{optim:StopTgtPos:Height:102.5}Electrons from 100 to 105 MeV/c]{%
\includegraphics[width=0.95\textwidth,trim=1.15cm 0.05cm 0.3cm 0.180cm,clip]{figs/optimisation/StopTgtPosition/WithBeamBlocker-Height-VaryShifts-Momentum_102-5.pdf}}%
\\\subfloat[][\figlabel{optim:StopTgtPos:Height:102.5-zoom}Electrons from 100 to 105 MeV/c (Zoomed)]{%
\includegraphics[width=0.95\textwidth,trim=1.15cm 0.05cm 0.3cm 0.180cm,clip]{figs/optimisation/StopTgtPosition/WithBeamBlocker-Height-VaryShifts-Momentum_102-5-Zoom.pdf}}%
\caption{\figlabel{optim:StopTgtPos:Height}
The effect of stopping target position on the height of electrons with a fixed momentum as they pass through the Electron Spectrometer.
The size of the variation indicates the stability of the dipole tune; the two parameters are clearly correlated.
Also striking---particularly in \protect\subref{fig:optim:StopTgtPos:Height:102.5-zoom}---is the way the dependence on the helical pitch angles is affected by the stopping target position.
}
\end{figure}
}

%\newcommand{\FigOptimStopTgtPosSensitivityIntegral}{
%\begin{figure}[tb]
%\centering 
%%	\fbox{
%\includegraphics[width=0.9\textwidth,trim=1.2cm 0.05cm 0.9cm 0.6cm,clip]{figs/optimisation/StopTgtPosition/Tidied_-AcceptedMomentum-Integrated_adjacent.pdf}
%%}
%\caption{\figlabel{optim:StopTgtPos:AcceptedIntegrated}
%The variation in sensitivity to different momentum electrons as a function of the target position, with respect to the nominal location.
%The darkest red line towards the top of the plot represents the signal sensitivity, and it is that line that should therefore be maximised.
%}
%\end{figure}
%}
%
%\newcommand{\FigOptimStopTgtPosSensitivityIntegralRatio}{
%\begin{figure}[tb]
%\centering 
%%	\fbox{
%\includegraphics[width=0.9\textwidth,trim=1.2cm 0.05cm 0.9cm 0.6cm,clip]{figs/optimisation/StopTgtPosition/Tidied_-AcceptedMomentum-Integrated_ratios.pdf}
%%}
%\caption{\figlabel{optim:StopTgtPos:AcceptedIntegratedRatio}
%The relative avvep
%}
%\end{figure}
%}

\newcommand{\FigOptimDIOBeamBlockGeometry}{
\begin{figure}[tb]
\centering 
%\fbox{
\includegraphics[width=0.8\textwidth]{figs/optimisation/BeamAndDIOBlocker/StopTgt_to_Detector-cropped.png}
%}
\caption{\figlabel{optim:DIOBeamBlock:Geometry}
Location of the beam blocker and one possible geometry for the \ac{DIO} blockers, both highlighted in green, shown here before optimisation.
}
\end{figure}
}

\newcommand{\FigOptimDIOBeamBlockESTDispersion}{
\begin{figure}[tb]
\centering 
%\fbox{
\includegraphics[width=0.8\textwidth,trim=0.3cm 0 1.7cm 0.6cm,clip]{figs/optimisation/BeamAndDIOBlocker/Tidied_ElectronDispersion-EST.pdf}
%}\\\fbox{
\includegraphics[width=0.8\textwidth,trim=0.3cm 0 1.7cm 0.6cm,clip]{figs/optimisation/BeamAndDIOBlocker/HandTidied_ElectronDispersion-EST-envelope.pdf}
%\includegraphics[width=0.8\textwidth,trim=1.5cm 0 8.5cm 3.0cm,clip]{figs/optimisation/BeamAndDIOBlocker/Tidied_ElectronDispersion-EST-envelope.png}
%}
\caption{\figlabel{optim:DIOBeamBlock:ESTDispersion}
Momentum-dependent dispersion of electrons passing through the spectrometer.
Top plot: the mean height of different momenta electrons as a function beamline distance, showing how the drift of the centre of gyration is truly proportional to the momentum.
Bottom plot: single standard deviation bands for electrons at different momentum, which shows how the envelope for different momenta overlap considerably, reducing the effectiveness of any collimators.
%Whilst the drift of the centre of gyration is proportional to momentum (top plot), there is reasonable overlap in the overall enevelope
%Mean height of electrons with different momentum from the stopping target to the detector before adding any DIO blockers and before optimised of the beam blocker.
}
\end{figure}
}

\newcommand{\FigOptimDIOBeamBlockAcceptances}{
\begin{figure}[tb]
\centering 
%\fbox{
\subfloat[\figlabel{optim:DIOBeamBlock:Acceptance:40}40 MeV/c]  {\includegraphics[width=0.35\textwidth,trim=0.0cm 0.3cm 0.5cm 1.5cm,clip]{figs/optimisation/BeamAndDIOBlocker/Tidied_Acceptance2D_40.pdf}}\hspace{1cm}
\subfloat[\figlabel{optim:DIOBeamBlock:Acceptance:60}60 MeV/c]  {\includegraphics[width=0.35\textwidth,trim=0.0cm 0.3cm 0.5cm 1.5cm,clip]{figs/optimisation/BeamAndDIOBlocker/Tidied_Acceptance2D_60.pdf}}\\
\subfloat[\figlabel{optim:DIOBeamBlock:Acceptance:80}80 MeV/c]  {\includegraphics[width=0.35\textwidth,trim=0.0cm 0.3cm 0.5cm 1.5cm,clip]{figs/optimisation/BeamAndDIOBlocker/Tidied_Acceptance2D_80.pdf}}\hspace{1cm}
\subfloat[\figlabel{optim:DIOBeamBlock:Acceptance:100}100 MeV/c]{\includegraphics[width=0.35\textwidth,trim=0.0cm 0.3cm 0.5cm 1.5cm,clip]{figs/optimisation/BeamAndDIOBlocker/Tidied_Acceptance2D_100.pdf}}
%}
\caption{\figlabel{optim:DIOBeamBlock:Acceptance}
Acceptance into the straw tracker for electrons with different momentum at the stopping target as a function of the beam and DIO blocker dimensions.
Note the logarthmic scale for the colour bar.
}
\end{figure}
}

\newcommand{\FigOptimDIOBeamBlockHitRate}{
\begin{figure}[tb]
\centering 
%\fbox{
%\hspace{-1.3cm}
\subfloat[\figlabel{optim:DIOBeamBlock:HitRate}Hit Rate from DIO]               {\includegraphics[width=0.45\textwidth,trim=0.8cm 0.3cm 0.3cm 1.05cm,clip]{figs/optimisation/BeamAndDIOBlocker/HitRate2d.pdf}}\hspace{0.04\textwidth}
\subfloat[\figlabel{optim:DIOBeamBlock:HitVAccept}Signal Acceptance vs Hit Rate]{\includegraphics[width=0.45\textwidth,trim=0.8cm 0.3cm 0.3cm 1.05cm,clip]{figs/optimisation/BeamAndDIOBlocker/HitRateVsAcceptance2dRatio.pdf}\vspace{1cm}{}}
%\subfloat[\figlabel{optim:DIOBeamBlock:HitVAccept}Signal Acceptance vs Hit Rate]{\includegraphics[width=0.45\textwidth,trim=0.8cm 0.3cm 2.5cm 1.05cm,clip]{figs/optimisation/BeamAndDIOBlocker/HitRateVsAcceptance2d.pdf}\vspace{1cm}{}}
%\hspace{0.12\textwidth}
%}
\caption{\figlabel{optim:DIOBeamBlock:HitRateAcceptance}
\protect\subref{fig:optim:DIOBeamBlock:HitRate} Number of straw tracker hits per DIO electron.
\protect\subref{fig:optim:DIOBeamBlock:HitVAccept} Ratio between the high-momentum electron ($p>100$~MeV/c) acceptance to the number of hits per DIO electron.  Colour is on a logarithmic scale.
}
\end{figure}
}

\chapter{COMET Phase-II: Optimisation}
\sectlabel{phaseII-optimisation}
%\begin{easylist}
%# Before a substantial sensitivity estimate can be made, need a solidly optimised design
%# Aiming for \sensePII within a single year of running
%# Designs previously optimised \cite{CDR}, and these results are used as nominal design / starting point
%# Fresh optimisation using new software / simulation, updated fieldmaps, physics lists and geometry
%# Some aspects fixed already since \phaseI under construction: Experiment hall, Torus1, detector solenoid, fieldmap and coil parameters?
%# Key areas for optimising:
%## Production target dimensions
%## Torus1 dipole field strength
%## Torus2 dipole field strength
%## Electron spectrometer dipole field strength
%## collimator shapes and locations
%## stopping target and beam blocker position and form
%## DIO blockers on spectrometer
%\end{easylist}
The last study into the sensitivity of \COMET \phaseII was performed in 2009~\cite{CDRphase2}, before the staged approach and \phaseI design had even been considered.
That study found that a \ac{ses} of $2.6\times10^{-17}$ could be achieved in $2\times10^{7}$ seconds of running, with a total expected background count rate of fewer than 0.34 events over the entire run period.
Since then, the collaboration's focus has shifted to design and R\&D for \phaseI and no further studies have been made of the \phaseII design.

The purpose of this chapter is to make use of the updates in the fieldmap calculation, the geometry handling, and physics modelling to revisit the design of \phaseII and demonstrate that it can do at least as well as the previous 2009 design.
In addition to updates in the simulation, some aspects of the actual design have been refined and fixed alongside \phaseI preparation, such as the experiment hall and the coils and cold-mass of the first stages of the muon beamline.
The aspects of the experiment that remain open for optimisation are shown in \tab{optimisation:possible-parameters}.

As an initial configuration, much of the design from the 2009 study will be used, with updates included for the areas of the experiment that have since been refined by the \phaseI design.
\TabOptimisationParameters%

\section{Optimisation Strategy}
%\begin{easylist}
%# Take some aspects as fixed
%# Limit scope and approach:
%## Ideally, each aspect optimised in combination to maximise signal acceptance and reduce background
%## How decoupled are each section?
%## In practise such an optimisation is not easy to do, instead aim to produce a baseline optimisation so that all backgrounds / issues can be identified
%## This can then form basis for further optimisation, with perhaps a smarter more integrated approach
%# Method:
%## Production target optimisation
%### Maximise muon and pion yield between 0 and 80 MeV at entrance to muon beamline
%### Parameters to vary: target length, target radius
%## Muon beam optimisation
%### Maximise muon stopping rate in stopping target
%### Minimise pion stopping rate
%### vary dipole along TS2 and TS4
%### vary Collimators: TS2 and at TS3
%## Electron spectrometer optimisation
%### Optimise dipole to increase signal acceptance
%### Optimse DIO blockers so DIO rate per straw is less than 1 kHz
%### Vary solenoidal field to increase separation?
%## Stopping target / beam blocker optimisation
%### Maximise reflection of signal electrons from upstream by tuning target position
%## Detector optimisation
%\end{easylist}
To perform a comprehensive optimisation of \phaseII, there are many parameters to be changed, corresponding to the shapes, positions, field strengths and so on of the different regions along the beamline.
\Tab{optimisation:possible-parameters} lists some 32 parameters, each one of which could be studied.
And yet, this number alone does not describe the full challenge of optimising the \phaseII design, since many of these parameters will correlate to one another.
For example, a correlation likely exists between the position and shapes of the muon beam collimator(s), the stopping target and the beam blocker, since all involve removing or stopping muons and other particles in the beam.
Other less intuitive correlations might also exist and a complete optimisation should be able to include the impact of these as well.

A full optimisation study then would involve a scan through a parameter space with at least 32 dimensions.
A brute force search of such a space would be nightmarishly slow and require an enormous amount of computing power.
Machine learning algorithms or intelligent scanning techniques might be able to tackle such a problem, and perhaps in the future these methods will be used.
In the meantime, however, we make use of the technique known as `physicists intuition' to approach the problem, whereby some parameters are assumed uncorrelated whilst others are disregarded on the expectation that their impact be small.
%This should be sufficient to establish a reasoned baseline geometry, whilst allowing for future work to refine the design further.
The goal of this chapter therefore is not only to optimise the experiment but also to evaluate the correlations.

This challenge is further complicated because it is not just the signal sensitivity that must be maximised, but the background rate that must be kept low.
When this work began, no start-to-finish simulation existed to estimate the signal and background rates.
As such, this study will progress by focussing on optimising the signal acceptance in the first instance, and then move to background estimation described in chapter~\sect{backgrounds}.

%Ideally this will also identify which parameters are the most strongly correlated, which would help make future optimisation studies more efficient.
The outputs of this optimisation should not be considered as final but should instead be treated as a baseline from which more intelligent approaches or physicists can iterate and improve.

%\section{Optimisation Goals}
%\begin{easylist}
%# Set sensitivity goal and optimise to reach this
%# Identify which parameters are strongly correlated to help focus and improve the efficiency of future optimisation studies.
%# Single event sensitivity only considers signal acceptance, but also need to understand backgrounds in terms of final confidence limit that can be set
%\end{easylist}

\section{Production Target Optimisation}
In the \phaseII \ac{CDR}, the production target is given as being 16~cm in
length and 4~mm in radius~\cite{CDRphase2}.  Since then, there have been changes to
the magnetic field in this region as well as the lengths and locations of
solenoids, shielding and beam-pipe, and the proton beam.  Previous studies have
looked at comparing the tungsten target proposed for \phaseII to other
materials~\cite{AEdmondsThesis}, and also drawn a comparison between
MARS~\cite{MARS1995}, Geant4~\cite{Geant42003} and the limited data available.

\subsection{Production Target Simulations} 
The goal of this study is to maximise the total muon and pion
yield below 80~MeV at the entrance to the Torus1 bent solenoid, by varying the
radius and length of the production target.
Whilst it is really the rate that negative muons stop in the Stopping Target that must be optimised, 
since it is known that muons above about 40 to 50~MeV/c will not stop, and since the bent transport solenoids themselves need optimising
we assume here that by maximising the number of muons and pions below 80~MeV/c will also maximise the stopping rate.
In addition this had the advantage of increasing the speed of the simulation.
The production target is one of the more computationally intensive aspects
of the experiment, given the relatively high energy protons and hadronic nature of their interactions with the target, which together result in a large number of secondaries.

The production target is formed from a solid cylinder of tungsten.
To find the optimal target geometry the radius and lengths were both varied individually, and then a combined scan performed where the length was held at the observed optimum. 
The location and orientation of the target were held fixed, since the axis of the proton beamline is fixed to intercept the muon beam axis
at a given point.  Once a more realistic proton beam becomes available, the location and direction of the production target
would also benefit from optimisation, however.  During the scan over
length, the back face of the target was kept 8~cm away from the muon beam since
the radiation shielding has previously been optimised.
%, and beyond this
%point the magnetic field will no longer be able to capture the pions and muons
%produced.

Primary protons were generated according to the description in the 2014 \phaseI TDR~\cite{TDR2014}.
This provides a double Gaussian position profile as well as the spread in momenta.
It does not discuss any dispersion or divergence in the beam, however.
The proton beam is currently under study, which is why the description at this stage is fairly simplistic.
The impact of a more realistic proton beam description will need to be studied in future work.

Protons in the simulation originated from a plane but, since there is some scope to tune the proton
beam's position, the input particle plane was moved to remain 1~cm away from
the front surface of the target.  Since the aim is to maximise the muon and
pion yield by varying only the length and radius, shifting the proton beam
input plane in this way removes any variation of target acceptance due to
divergences of the proton beam in the field before the target.

\FigOptimProdTgtLength
\subsection{Length Scan}
Different length targets were simulated with \num{3e5}~\ac{POT} per length.
The target length was varied in steps of 4~cm from 4 to 32~cm, whilst the target radius was held fixed at the CDR value of 4~mm.

\Fig{optimisation:ProdTgtSec:Length:Momentum} shows the momentum distributions of pions and muons for different target lengths, which are 
given as half-length following the Geant4 convention.  
\Fig{optimisation:ProdTgtSec:Length:Integral} then shows these distributions integrated up to different momenta.
From these plots it can be seen that for both muons and pions, the optimum target length occurs around a total length of 32~cm.

Additionally it can be seen from
\fig{optimisation:ProdTgtSec:Length:IntegralRatio} that the shape of the
momentum distributions changes only weakly as a function of the target length.
These plots were produced by normalising the integrated momentum contours of
\fig{optimisation:ProdTgtSec:Length:Integral} to the total integral below
400~MeV.  As a result, it is possible that the actual shape variation is even
weaker than apparent here, since in the present sample size, the high momentum
tail is not well sampled at small target lengths, such that a skew in the
normalisation might occur.

\FigOptimProdTgtRad
\FigOptimProdTgtFinal
\subsection{Radius scan}
In parallel to the length optimisation scan, different radii targets were also simulated.
Targets with radii of 2, 4, 6, 8, 10, 14, 18, 22, 26, and 30~mm were tested.
The target length was held at the \ac{CDR} value of 16~cm in total.

The results of these scans can be seen in \fig{optimisation:ProdTgtSec:Radius:Momentum} and \fig{optimisation:ProdTgtSec:Radius:Integral},
where it can be seen that a maximum in both the muon and pion yields at the entrance to the Torus1 section is achieved at a radius of about 10~mm.
As in the length scan, the varation in the shape of the momentum distributions is rather weak as a function of target radius.

\subsection{Final Result}
\FigOptimProdTgtComparePhases
Since the length and radius scan were performed in parallel, a final cross check was performed where the optimal radius was confirmed at the optimised target length.
The integrated spectrum is shown in \fig{optimisation:ProdTgtSec:Final:Integral} where it can be seen that the optimum radius once the target length is increased to 32~cm is still 10~mm.

\Fig{optimisation:ProdTgtSec:Phase1vs2} shows the total muon and pion yields at the entrance to the Torus1 for the final optimised \phaseII target and compares this to the optimised \phaseI target design.
The increased yield in the low momentum range is due to the heavier target nucleus which produces more low-momentum pions in the backwards direction.
Since only muons below around 50~MeV/c tend to stop in the target, this increase in the low momentum yield amounts to a factor 2 or so gain in the stopping muon rate per POT.

\section{Dipole Strengths of the Muon Beamline}
%\begin{easylist}
%# Production target simulation
%	## Plots of input distribtion (possibly shown in previous section)
%	## Mention gas pressure at exit of ProdTgtSec but show why this shouldn't have been an issue
%# Optimisation technique
%	## Vary dipole field strength in Torus1 and Torus2 simultaneously
%	## Mesh scan through parameter space
%	## Look at muon stopping rate
%	## Stopping target as in the CDR 
%	## No collimator material included
%# Results
%	## Muon stopping rate vs dipole field plot
%	## Pion stopping rate vs dipole field plot
%	## Discussion of corraltion
%# Selected optimal value
%\end{easylist}

The full 180\degree bent solenoid that makes up the bulk of the muon transport beam line is actually broken down into two
90\degree pieces---Torus1 and Torus2 (also known as TS2 and TS4 by the magnet group).
Each of these sections has its own dipole field which need not provide the same dipole field strength.
The Torus1 section has been built already and was previously optimised for both \phaseI~\cite{TDR2016} and \phaseII~\cite{CDRphase2}. 
The dipole coils that it contains are designed to produce a dipole field of 0.055~T.
By running a lower current through these coils one could reduce the Torus1 dipole field without too much effort. 
However if \phaseII should require a greater dipole field strength for this region
that might be trickier since it could require additional windings to be inserted---this would not be impossible but might costly.

\subsection{Large-sample Production Target Simulation}
\sectlabel{largeSampleProdTgt}
Before the muon beam section could be optimised, a large set of \ac{POT} events were produced transporting all particles up to the entrance of the muon beam section.
In total, \num{2.3e8}~POT were simulated through this stage, equivalent to about 1.5 bunches at \phaseII.
The production target used the optimised geometry from the previous section.
All particles that hit the surface of the Torus1 container volume were read out for later re-use in a way that ensured double-counting of particles could not occur.
In addition, particles entering the proton beam dump were also saved if later simulations wished to study their impacts.

\subsection{The Optimised Dipole Field Strengths}
The figure-of-merit for this optimisation is the muon stopping rate, which will be greatest for the optimal field configuration.
%Toshiba have previously performed a realistic field calculation using OPERA for a single coil whose on-axis field has a mean of about 0.055~T.
To identify such a configuration, a two-dimensional scan over the Torus1 and Torus2 dipole field strengths was performed.
To vary the field strength a scale factor was applied, for each of the 90\degree section, to the realistic dipole field description that was provided by Toshiba and described in \sect{sw:fieldmap}.
These scale factors were varied in steps of 0.125 (equivalent to 6.875~mT) and for each combination muons and pions from $8\times10^6$~POT were transported to the stopping target.
%The stopping rate for each combination of scale factors was then assessed.

\FigOptimMuBeamDipoleMuStops
\FigOptimMuBeamDipolePiStops

The muon stopping rate as a function of the two dipole field strengths is shown in \fig{optim:muBeamDipole:stoppedMu}.
The horizontal axis in that figure shows the scale factor applied to the Torus1, the 90\degree bent solenoid already built for \phaseI, whilst the vertical axis shows the scale factor for the Torus2.
A scale factor of 1.0 means the scale factor will be same as the current \phaseI design.
Since there is little improvement by moving to larger values for Torus1 the optimal scale factors are chosen to be 1 and 0.5 for the Torus1 and TTorus2 respectively.  
This translates to dipole field strengths of 0.055 and 0.0275~T respectively.
%\CHECK{Are these the final selected dipole field strengths?}

Also striking from \fig{optim:muBeamDipole:stoppedMu} is the anti-correlation between the two dipole field strengths.  
Roughly speaking the sum of the optimal dipole field values is constant, \ie $B_1^\text{optimal} + B_2^\text{optimal} = \text{const}$.
This correlation had not been seen before, perhaps due to the lack of computing power necessary to scan such a large phase space.
Such a correlation can be understood in the following way:
the geometry of the stopping target geometry (and its position) is kept fixed during the dipole scan.
This essentially fixes the upper momentum of muons that can stop in the target.
%The initial momentum distribution grows roughly linearly with momentum for the momenta of muons that will stop in the target (see the later sections for more on this), which would suggest tuning to the higher momentum would increase the number of muons reaching the target.
%That effect however is greatly reduced by the fact that the initial muon beam has a very large transverse position profile
The drift of a particle due to the dipole field is proportional only to the distance travelled in that dipole field and the dipole field strength, and does not depend on the particle's momentum.
This means that the total drift due to the Torus1 and Torus2 dipole components is proportional to the weighted sum of the two dipole field strengths, where the weight is the distance travelled through each of the two sections.
Since each section is the same length, the time for a particle to pass through Torus1 will be roughly the same as the time in Torus2 so that the weights would be roughly equal.
This leads to the total drift being proportional to the sum of the two dipole fields.
Since the production target (as the source of muons) and stopping target (as the muons' destination) are at a fixed height with respect to one another the optimal vertical drift is also fixed so that the sum of the dipole field strengths should also be roughly fixed.

With that said, the Torus1 section has a higher pion flux which causes some asymmetry between the two sections.
Keeping more pions on-axis in the Torus1 section means that more muons will enter the Torus2 section from those pions that have decayed.
But since the pion momentum distribution is slightly higher than the muon distribution keeping pions on axis requires a larger dipole field strength.
This could explain the slight asymmetry where the muon stopping rate appears slightly larger if Torus1's field is larger than that of Torus2.

It is also interesting to consider the pion stopping rate as a function of the dipole field strengths.
However, as can been from \fig{optim:muBeamDipole:stoppedPi} the stopping rate is close to the level of POT events used in the simulation so that the plot is dominated by statistical fluctuations.

\section{Electron Spectrometer's Dipole}
%\begin{easylist}
%# Method
%	## Realistic muon stopping distribution
%	## Signal electrons at the target 
%	## vary the dipole field strength and study the acceptance
%	## Uniform dipole field with sharp turn-on
%# Signal height vs. dipole field strength
%	## 2D plots for no dipole, 0.1 T and 0.2 T
%	## Stack of 1D mean heights
%# Signal acceptance vs. Dipole field
%	## Survival probability vs beamline pos
%	## Overall acceptance at the detector
%# Mean height for other momenta
%\end{easylist}
The next element in the beamline after the muon transport solenoids will be the stopping target.
However, in order to study the impact of changing the stopping target parameters one will need to look at the impact on the signal acceptance into the detector.
To study that requires the components of the beamline intermediate to the target and the detector be optimised, namely the electron spectrometer.
The key free parameter in this section is the dipole field strength along the spectrometer.
The solenoidal field and solenoid aperture could also be optimised in principle, but these are considered beyond the scope of this study at this point and so the CDR values are held fixed here.
The point of this section is to establish the optimal dipole field strength given fixed target parameters which will then be studied separately together with the stability of the dipole field tune checked.

\subsection{Method and Potential Short-comings}
To study the effect of the dipole field on signal acceptance, a realistic muon stopping distribution in the target was produced by transporting muons from the production target simulation through to the stopping target.
Signal electrons were then injected at the target with the realistic stopping distribution and propagated through the beamline to the detector with different dipole field strengths.

A non-trivial short-coming of the current study is that the dipole field along the spectrometer is poorly modelled---no realistic coil simulation exists, unlike for the bent muon transport beamlines.
As a result, a perfectly uniform dipole field is assumed with a sharp switch on and off at the entrance and exit of the spectrometer.
The impact that this has on the final result is not clear: one might expect it to be small given the relative strengths of the dipole and solenoidal fields and overall it is the integrated field that tends to matter.
However, the sharp switch-on of the field at the entrance and exit of the spectrometer is clearly not physical.
Given that the gradient introduced in the field by bending is present before the actual entrance and exit of the spectrometer (as a fringe field), some drift can be anticipated in this region.
A realistic dipole field with a realistic fringe field might overcome some of this drift however, so that the uniform field used here will most likely not capture this effect.
Nevertheless, given the absence of a realistic dipole field and scope of this study, the uniform one is the only real option at this point.

\subsection{Results}
\FigOptimESTDipoleBeamHeightTwoD
\FigOptimESTDipoleBeamHeightMean
\Fig{optim:ESTDipole:Beam} shows the projection of electron trajectories to the beam axis coordinate system for three different dipole field values.
The potency of this approach is clear from these plots; the tuneable dipole fields allow the momentum of electrons which remain on-axis to be accurately controlled (during run-time), which will benefit systematic and calibration studies that wish to study the \ac{DIO} spectrum at a lower energy.
\Fig{optim:ESTDipole:MeanHeight} then collects these plots with other dipole field strengths, plotting the mean height for all simulated electrons against the distance along the beam axis.
From this plot it can be seen how a dipole field of about 0.18~T appears optimal to keep the signal electrons on axis.

The probability for electrons to reach a given point in the beamline is shown in \fig{optim:ESTDipole:MeanFlux}, and indeed from this it can be seen that to maximize the probability of an electron reaching the detector
a dipole of around 0.18~T is desirable.
The behaviour of the low dipole field values (0 to 0.08~T) in this plot was not expected, but it is believed this is an artifact of the way this plot is made, coupled with a degree of mirroring at the entrance to the spectrometer which is enhanced as the dipole field strength increases.
If correct, a realistic dipole field calculation would be important to quantify and confirm this behaviour.
\FigOptimESTDipoleBeamFluxMean

\FigOptimESTDipoleAcceptanceVsDipole
Finally, to confirm the optimal dipole field strength the true geometric acceptance of the detector system is checked as a function of the dipole field strength, which is shown in \fig{optim:ESTDipole:acceptance}.
An electron is considered to have been geometrically accepted by the detector in this simulation if it produces at least one hit in the detector system.
In principle this could be in any straw plane, but in practice this is almost always in the first layer of straws.
Since this is a different way to analyze the acceptance compared to the survival probability it would not suffer to the artifact seen for low dipole field values in \fig{optim:ESTDipole:MeanFlux}.
Nonetheless, \fig{optim:ESTDipole:acceptance} confirms that the optimal dipole field strength is very close to 0.18~T.

A second important conclusion can be drawn from the fact that the dependence on the dipole field strength is relatively weak around the optimal value of 0.18~T.
A change of about 10\% in the dipole field strength only reduces the signal acceptance by about 3\% whilst a change of about 5\% would see a reduction of only about 0.7\%.


\section{Stopping Target Position}
\sectlabel{optim:stopTgtPos}
%\begin{easylist}
%# Method
%	## Shift target position up or downstream by 50 cm 
%	## Muon stopping distribution with realistic beam sim. only include muons for speed
%	## Electron acceptance using realistic stopping distribution
%	## Don't change anything else about the target (size, length, number of disks, beam blocker distance, etc)
%	### Too many parameters, want only to achieve a complete baseline here.
%# Muon stopping rate vs, position
%# Electron acceptance times muon stopping rate = signal sensitivity vs. position
%# Stability of dipole field tune 
%# Impact on low energy electrons with regards to the subsequent DIO blocker tune stability
%# Impact of beam blocker on acceptance
%## See appendix for further discussion and further optimisation in this region
%\end{easylist}
The final aspect to be studied with regards to maximizing the signal sensitivity is the stopping target.
In principle there are many parameters related to the stopping target such as the location, disk shape (profile and thickness), and disk spacing
In addition the beam blocker should be considered in parallel with the stopping target since it sits so close to the target itself and can be expected to have a big impact on the signal acceptance.
However, this leaves far too many parameters to be considered all at once.

Since the field around the target tapers sharply various competing factors must be considered.
For example, prior to the stopping target region the muon beam is transported through a 3~T solenoidal field.
The magnetic field in the target region, however, tapers to about 1~T, which would cause the envelope of the muon beam to grow.
Moving the stopping target downstream would mean that the muon beam arrives with a larger aperture, and would therefore prefer a larger stopping target or else fewer muons will actually hit the target.
On the other hand, from the perspective of signal acceptance, the tapered field can be used to mirror signal electrons that are initially produced heading upstream therefore increasing the signal acceptance.
Moving the target further upstream then will reduce this effect as the difference between the magnetic field strengths at the exit of the bent muon transport solenoid and at the stopping target is reduced.

Given this and the need to reduce the number of parameters inspected, the target and beam blocker design was held fixed in this study and only the position was changed by moving the target upstream and downstream by $\pm50$~cm with respect to the nominal target location as given in the CDR~\cite{CDRphase2}.
Given that the target disks will occupy in total about 1~m, a shift of 50~cm corresponds to half the target length.
In each different position of the stopping target, as for the spectrometer dipole optimisation, a realistic stopping distribution was produced by running muons from the large-scale production target simulation through to the target.
This stopping distribution was then re-used to introduce signal electrons accordingly.
Additionally, low momentum electrons were also studied in order to study the impact of target position on the height of both signal and low-energy electrons as they pass through the spectrometer.
This is important both to check the correlation of the dipole field tune with the stopping target position, but also how the subsequent DIO blocker height optimisation will correlate the target position.
\FigOptimStopTgtPosMuStops

\Fig{optim:StopTgtPos:MuStops} shows how the rate of muon stops per \ac{POT} is affected by changing the position of the stopping target.
The relationship is roughly linear, dropping from around $2.2\times10{-3}$ muon stops per POT when the target is shifted upstream by 50~cm to about $1.3\times10^{-3}$ muon stops per POT if the target is shifted 50~cm in the other direction.
This relationship is as expected given the fixed radius of the target and the growth of the muon beam aperture arising from the reduction in the field strength.

\FigOptimStopTgtPosSensitivitySpect
In \fig{optim:StopTgtPos:AcceptedMomSpect} one can see the way the electron acceptance changes for different target positions.
Acceptance here is defined as producing at least one hit in the detector and the momentum shown is the momentum at the target, which is not necessarily the same as the momentum at which the electron is observed.
Each histogram in \fig{optim:StopTgtPos:AcceptedMomSpect} is normalised to the number of primary electrons introduced at the target per MeV/c and then scaled to the muon stopping rate.
This normalisation makes the value of each curve proportional to the sensitivity of the experiment to different momentum electrons, up to factors such as analysis cuts like timing and reconstruction quality.

\FigOptimStopTgtPosSensitivityIntegral
Since the parameter we wish to optimise here is the location of the stopping target, \fig{optim:StopTgtPos:AcceptIntegral:Adjacent} represents the same data as in \fig{optim:StopTgtPos:AcceptedMomSpect} but with each line representing the content of a different 5~MeV/c bin as a function of the target position.
For signal, it is the 105~MeV/c line (dark burgundy) that is most important and it can be seen that this is optimised for shifts upstream of the nominal position from between 10 and 20~cm.
It is also interesting to note that the acceptance of lower energy electrons is relatively decreased as the target is moved upstream, as can be seen in \fig{optim:StopTgtPos:AcceptIntegral:Ratio}.
This could be useful as a way to suppress hit rate from DIO electrons later.

Whilst we do not intend to optimise the beam blocker at this point, to check the impact that it has on the optimisation of the target position simulations were performed where the blocker was completely removed.
\Fig{optim:StopTgtPos:AcceptedMomSpectNoBeamBlock} shows the product of the stopping rate and electron acceptance when the blocker is removed.
From this the trend is much cleaner, electrons below 70~MeV/c are suppressed as you move the target upstream whereas the high energy electron acceptance increased.
\FigOptimStopTgtPosSensitivityNoBeamBlock


\FigOptimStopTgtPosHeights
Finally the relationship between the stopping target position and the mean height of electrons through the spectrometer is demonstrated in \fig{optim:StopTgtPos:Height}.
Each plot shows the mean heights for electrons with a given momentum for different stopping target positions and it is clear that there is some correlation between the mean height and the position of the stopping target.
For this reason, a more complete optimisation should consider optimising the two parameters simultaneously.
Nonetheless the change is not particularly large, only about a few cm difference by te end of the spectrometer for signal electrons.
This correlation is also likely related to the way the acceptance changes for different target positions when the beam blocker is removed.

The second striking feature from the plots in \fig{optim:StopTgtPos:Height} is the way the mean height aquires a strong sinusoidal component for large target shifts upstream.
This suggests that when the target is shifted upstream the electrons passing the entrance to the spectrometer tend to have a particular value for the pitch and phase angles of the trajectories.
Several separate mechanisms could produce this effect.
Firstly the acceptance around the target itself could aquire a stronger preference for certain pitch and phase angles when the target is moved upstream.
Secondly, since the stopping target disks will see more of the muon beam upstream, the muon stopping distribution could become less homogenous.  Whilst electrons are produced isotropically path with less target material along them will tend to accept outgoing electrons more readily such that if more muons stop at one side of the target than the other a dependence on pitch and phase could arise.

Clearly then the stopping target region is a very complicated area; even though this study has focussed on a single parameter---the position of the target itself---it is clear that this correlates to many other variables.
This is a region in the experiment most ripe for further optimisation then, and indeed appendix~\ref{sec:appendix:stopTgtImprove} shows some of the first steps that have made in this direction since the optimisation described here was completed.
Unfortunately, at the time this work was carried out an error in the normalisation of these plots lead to the conclusion that the optimal shift was between 0 and 10~cm upstream.
Given the complexity of the optimisation in this region, it was therefore decided to keep the stopping target at the nominal location for the subsequent steps.
Having corrected the normalisation of the plots, the conclusion now is that the optimal location is between 10 and 20~cm upstream, so perhaps the target should have been shifted back.
However the improvement to the sensitivity would have been small: shifting the target back about 20~cm improves the sensitivity by around 2\% compared to the signal acceptance at nominal position.
%Given this small improvement and the complexity of this region it was decided that the stopping target was kept to the nominal position.

\section{Collimators in the Muon Beamline}
\FigOptimMuBeamDipoleMuDispersion
%\begin{easylist}
%# Muon beam without collimators
%	## Height vs momentum before and after 180 degrees
%	## Beamline projection for stopped and high-P muons
%	## Subtraction of the good and bad muons to suggest the region of interest for collimators
%	## Why the asymmetry upstream?
%# Collimator optimisation
%	## Analytic collimator study to reduce simulation time
%	### Cannot study secondary particles
%	### Assumes all particles entering the collimators are killed
%	## Current collimator from Phase-I are in a different location to those suggested here
%# Results of study, recommended collimator design
%\end{easylist}
With the beam line optimised for high signal efficiency, one can look at improving the background by adding collimators into the muon beam line to reduce the flux of high-momentum muons and pions.

\Fig{optim:MuBeamDipole:MuDispersion} shows the dispersive effect of the bent solenoid and dipole field on muons passing through the beamline.
Thanks to this dispersion low momentum muons that can stop in the target and high momentum muons which could produce backgrounds are separated sufficiently so that material in the beam pipe can selectively remove the dangerous, high-momentum muons with only a small impact on the muon stopping rate.

\subsection{Collimator Placement}
\FigOptimMuBeamCollimMuonPaths
The plots in \fig{optim:MuBeamCollim:Beamline} give a sense of where best to locate the collimating material.
In \fig{optim:MuBeamCollim:Beamline:All} the paths of all muons along the beamline is shown.
\Fig{optim:MuBeamCollim:Beamline:Stopped} then separates out the muons that stop in the target which should be compared to 
\fig{optim:MuBeamCollim:Beamline:HighP} showing the paths of muons that reach the stopping target region with momentum greater than 70~MeV/c (the threshold for a muon to decay to an electron with $p>100$~MeV/c).

It is interesting to note the apparent asymmetry in the high-momentum muons at the entrance to the Torus1 that can be seen in \fig{optim:MuBeamCollim:Beamline:HighP}.  
This is not because of some momentum-dependence in the acceptance of the preceding beamline, but due to the fact that muons in that plot are only included if they are `dangerous' in the region around the stopping target.
Even without additional collimators, the beam pipe itself removes high momentum muons that enter in the lower half of the beamline.
The validity of only tagging high momentum muons around the stopping target region comes from the assumption that the products of high-momentum muons that decay before this region can be reliably removed.
It is important then that this assumption be checked, but for this work this is left as a task for the future.

Finally then in \fig{optim:MuBeamCollim:Beamline:Diff} the difference between the high-momentum muons and the paths of those that stop is shown. 
Regions in blue on this plot show where many more high-momentum muons pass than stopping muons; it is in these locations that collimator material will be most effective.
This approach could also be improved, since taking the straight difference between the two plots implies equal weighting for stopping muons and high-momentum ones.
In reality, whilst a muon tagged as stopping is definitely going to stop, a high-momentum muon should be weighted by the probability to produce a signal-like electron and the probability that this electron survives to be accepted into the detector, passing all analysis cuts.
This would make the weighting for the high-momentum muons be a function of the beamline distance itself which again requires a study into how high-momentum electrons are accepted.
However for the purposes of obtaining a qualitative sense of where to collimate the unweighted difference should be sufficient.

Two regions of interest appear: in the upper half of the entrance to Torus1 and the lower half of the exit of Torus2.  
Collimating at the Torus1 entrance is justified by the basis that high momentum muons will tend to have larger gyroradii compared to the muons that stop, and that at this point the beam is largely on-axis and has not yet been dispersed.
Collimating at the exit of Torus2 is readily understood on the grounds that the high momentum muons will have all drifted downwards by this point, compared to the low momentum stopping muons which are kept on-axis by the dipole fields.

\FigOptimMuBeamCollimTransverseSep
These conclusions are backed by the plots shown in \fig{optim:MuBeamCollim:TransverseSep} which show transverse slices through the beamline at the entrance to Torus1 (0\degree of bent solenoid), the midpoint between Torus1 and Torus2 (after 90\degree of bent solenoid) and at the exit of Torus2 (after 180\degree).
From these plots one can also see how in the middle of the bent solenoids (at 90\degree) the separation between muons that will stop and those that will have momentum greater than or equal to 70~MeV/c in the stopping target region is weakest, and hence collimators in this location will not be so effective.

\FigOptimMuBeamCollimTorusOne
\FigOptimMuBeamCollimTorusTwoFraction
\subsection{Collimator Height Optimisation}
To identify the optimum height for the collimators in a computationally efficient way, events were generated without any collimators included.
The full three dimensional trajectory of all particles as well as the decay tree were persisted.
This allows for a `virtual' collimator to be used, where particles that enter a defined region and their secondaries are removed from the downstream plots.
Whilst this method allows for many collimator shapes and heights to be tested quickly, it comes with a couple of limitations.
Firstly, the accuracy depends on the trajectory sampling density, which should be fine but this results in large data sizes.
Secondly, realistic material effects of the collimator cannot be captured, such as the probability a particle is simply scattered rather than stopping completely, or the result of secondary particles produced in the collimator itself.

\Fig{optim:MuBeamCollim:Torus1} shows the results of lowering the bottom edge of the collimator material in Torus1.
\Fig{optim:MuBeamCollim:Torus1:perPOT} shows the probability per \ac{POT} that different types of particle (stopped muons and high momentum muons, pions, and electrons) pass the collimator as a function of the collimator height.
On the other hand, \fig{optim:MuBeamCollim:Torus1:fraction} shows the same plots but normalised to the total number of each particle type that reaches the collimator in the first place, therefore showing the survival probability along just the collimator region.
Based on these plots, for a collimator that starts at 120~mm above the beamline axis, 14\% of high momentum pions are removed, high momentum muons are suppressed by 14\% whilst the muon stopping rate is reduced only by 3\%.

\FigOptimMuBeamCollimTorusTwoContours

For the second collimator at the exit to Torus2, the situation is slightly more complicated since in principle the optimum height could be correlated to the height of the upstream Torus1 collimator.
To account for this, the virtual collimator technique was applied for both the Torus1 and Torus2 collimator sections simultaneously.
As a result, the 1-dimensional plots of \fig{optim:MuBeamCollim:Torus1} become 2D as can be seen in 
%\fig{optim:MuBeamCollim:Torus2:perPOT}, which is normalised to the rate per POT, and 
\fig{optim:MuBeamCollim:Torus2:fraction}, which is normalised to the particle flux just before the collimator similar to \fig{optim:MuBeamCollim:Torus1:fraction}.
%\FigOptimMuBeamCollimTorusTwoPerPOT

\FigOptimMuBeamCollimMuonPathsWColl
\Fig{optim:MuBeamCollim:Torus2:contours} represents the stopping and high-momentum plots in a way that is easier to compare the two directly.
Each line in that plot is a contour showing a change of 2.5 percentage points to the yield.  Total acceptance, or 100\% is in the bottom right corner.
From this plot it can be seen that whilst keeping more than 97.5\% of the muon stopping rate (the bottom-right most blue contour), the maximum high momentum muon suppression is achieved when the Torus1 collimator sits about 140~mm above the beam axis, and the Torus2 collimator sits about 120~mm below it.
To be precise, at these collimator values  the muon stopping rate is kept at 99\% of the no-collimator rate, whilst high momentum pions, muons and electrons drop to 27.6\%, 20.9\% and 11\% (although this last value is very statistically limited) respectively, compared to their no-collimator rates.
At this point the selected values for the collimator heights give conservative background suppressions and tighter values could be chosen.
Given that backgrounds from high-momentum muons are suppressed compared to the actual rate of `dangerous' muons by the geometric acceptance of the remaining beamline and the timing and momentum cuts this seems reasonable at this stage although this will of course be investigated in the next chapter.

Finally, for comparison to the original plots, \fig{optim:MuBeamCollim:BeamWColl} shows the impact the new collimators have on the muon components of the beam.
It is clear how greatly reduced the number of muons passing the stopping target region has now become.
\clearpage

\section{The Beam and Decay-in-Orbit Blockers}
\sectlabel{optim:BeamAndDIOBlocker}
%\begin{easylist}
%# Dispersion of electrons through the spectrometer
%# Width of beam plot (1 sigma band around mean height for different momentum?)
%# Method for optimisation
%# Goal: Hit rate suppression, not backgrounds
%# Signal, DIO acceptance
%# Hit rate
%\end{easylist}
\FigOptimDIOBeamBlockGeometry
\FigOptimDIOBeamBlockESTDispersion
For every muonic atom formed in the stopping target some 39\% will undergo \acf{DIO}.
With some $1.4\times10^{8}$~\ac{POT} per bunch and a muon stopping efficiency around $1.7\times10^{-3}$ stops per \ac{POT}, one can expect about $9\times10^{4}$ DIO events per bunch.
If the detector system were exposed to this level of hits it would be impossible to resolve any signal electrons.
However, above the end-point energy for electrons coming from free muon decay, the DIO spectrum falls extremely quickly, with around 99\% of DIO electrons being produced with less than 59~MeV/c.
To this end the electron spectrometer's primary purpose is to disperse away the electrons below 60~MeV/c whilst keep signal electrons on-axis, such that material in the beamline can be tuned to remove the low energy DIO electrons.
The beam and DIO blockers are highlighted in green in \fig{optim:DIOBeamBlock:Geometry}.

\Fig{optim:DIOBeamBlock:ESTDispersion} demonstrates the dispersion that appears at the end of the spectrometer with the nominal beam blocker design and no DIO blockers in place.
The DIO blockers will be inserted along the bottom of the spectrometer and tuned to scrape away a sufficient number of DIO electrons, which will tend to travel towards the bottom of the beam pipe.
However whilst the centre of gyration drifts with a well-behaved proportionality in a bent solenoid, the actual separation between signal and \ac{DIO} electrons is in reality less clear as can be seen by the lower plot in \fig{optim:DIOBeamBlock:ESTDispersion}.

In addition to the DIO blockers along the spectrometer, the material of the beam blocker immediately after the stopping target also plays a role in suppressing the \ac{DIO} rate since low energy electrons remain closer to the beam axis, as can be seen in the lower plot of \fig{optim:DIOBeamBlock:ESTDispersion} around the stopping target.

This section therefore describes a simultaneous optimisation of the beam blocker radius and the DIO blocker height.
The overall goal is to suppress the DIO rate to less than a single DIO electron per bunch whilst maintain maximal signal acceptance.
As for the muon beam collimators, we use the analysis-based collimator approach, where no blocking material is included during simulation.
Instead a high-sampling density for particle trajectories is used such that during analysis particles and their secondaries can be `killed' if they enter a region that would contain material of either the DIO or beam blocker.

\FigOptimDIOBeamBlockAcceptances
\FigOptimDIOBeamBlockHitRate
The results of this study are shown in \fig{optim:DIOBeamBlock:Acceptance}, where the geometric acceptance into the detector is shown for four different electron momenta as a function of the beam blocker radius and DIO blocker height.
It is clear that for all values of the blockers' dimensions, electrons above 100~MeV/c have a much better acceptance.

The mean hit rate per DIO event is shown in \fig{optim:DIOBeamBlock:HitRate}.  
It is formed by, for each combination of DIO and beam blocker dimension, the weighted integral of the acceptance of electrons as a function of momentum with the mean DIO rate in each momentum bin.
To compare the hit rate to the acceptance of signal, the ratio between the hit rate and the high-momentum electron acceptance is shown in \fig{optim:DIOBeamBlock:HitVAccept}.
It is clear from this figure how quickly (note the logarithmic colour scale) the DIO hit rate is suppressed by increasing the DIO and beam blocker dimensions compared to the signal acceptance.

With a beam blocker of 24~cm radius and DIO blockers set to 35~cm below the beam axis, the DIO hit rate is about $2.2\times10^{-5}$ per DIO event, or about 2 DIO hits per bunch.
For the same blocker dimensions, the geometric signal acceptance is about 0.22\%.
Given the steep drop-off in hit rate versus signal acceptance around these values a finer scan in this region is an important check for the future.
What hit rate in the straw tracker is tolerable is a number for future studies.  In \phaseI, the straw tracker is expected to operate with a hit rate around 1~kHz per straw.
\phaseII will likely use finer straws, but with the current \phaseI design, some 133 straws occupy a layer, so that a total hit rate into the Straw Tracker of around 200~kHz should be acceptable.

Given the scope for future improvements in the granularity of the \phaseII straw tracker and the sensitivity to the beam and DIO blocker dimensions, the somewhat non-conservative values described above of 24~cm and 35~cm for beam blocker radius and DIO~blocker depth below the beam axis are selected here.

%\section{Revisiting the Torus2 Collimator after Beam Background Studies}
\section{Summary of optimised parameters}
The complete set of optimised parameters is shown in \tab{optim:AllParameters}.

\begin{table}[btp]
	\centering
	%\begin{tabular}{r|lp{6cm}}
		\begin{tabular}{lL{0.2\textwidth}L{0.35\textwidth}}
		Parameter & Value & Comments \\
		\hline
		\hline
		Production Target Length & 32~cm & \multirow{2}{4cm}{Placed asymmetrically about muon beamline axis} \\
		Production Target Radius & 1~cm & \\[1ex]
		\hline
		Torus1 (TS2) Dipole Strength & 0.055~T & same as \phaseI \\
		Torus2 (TS4) Dipole Strength & 0.0275~T & same as \phaseI \\
		\hline
		Torus1 Entrance Collimator Height& 14~cm above beam axis & From top of beam pipe downwards  \\
		Torus2 Exit Collimator Height& 12~cm below beam axis & From bottom of beam pipe upwards\\
		\hline
		Stopping Target Shift& 0~cm& Unchanged from CDR value (2.8~m from centre of final coil in Torus2 to front of target )  \\
		\hline
		Electron Spectrometer Dipole & -0.18~T& Negative compared to Torus1 and Torus2 dipole fields \\
		\hline
		Beam Blocker Radius & 24~cm &  \\
		DIO Blocker Height & 35~cm below beam axis& From bottom of spectrometer upwards  \\
		\hline
		\hline
	\end{tabular}
\caption{\tablabel{optim:AllParameters}
Optimised values for the parameters studied in this chapter.
Many more parameters remain to be optimised that were considered beyond the scope of the present work.
}
\end{table}

\section{Future optimisations}
The primary goal of this work is to update the previous optimisation from the 2009 CDR and provide a new baseline design.
However whilst touching every aspect of the layout of \phaseII, the optimisation developed here is not exhaustive and there is much scope for further work.

The following list is a short a summary of some of the areas that could be developed further:
\begin{description}
\item[ Refine the optimisation criteria]
	In many of these studies only the signal efficiency, or even some proxy, is used to identify the optimal value of the parameter under question.
	In reality one ought also to study simultaneously the background rates.  
	These two quantities should be combined into the expected confidence limits given a null or background-only measurement.
	At the same time other quantities such as cost and run-time will also need considering, although the latter is likely reduced simultaneously with maximising the overall sensitivity.
\item[ High energy electron acceptance vs. beamline distance ]
	When tuning the muon beam collimators, the key goal is to remove high-energy particles that can produce high energy electrons. 
	A particularly useful study would be to evaluate the acceptance to signal-like electrons which originate along the beamline.
	This should include those electrons which originate with momentum greater than 105~MeV/c but which arrive at the detector with signal-like energies.
	With this information, it becomes easier to identify how soon along the beamline one must collimate away the high-energy particles which may lead to improvements in the muon stopping rate or the signal acceptance.
\item[ Stopping target shape, thickness, and disk spacing ]
	This study has focussed only on the position of the stopping target.
	Clearly the actual shape should be studied as well. In particular, given the dispersion in the muon beam, and the changing solenoidal field strength around the target, a target design that uses disks of varying profile or even thickness has potential to improve the experimental sensitivity.  
\item[ Correlation between dipole fields and stopping target shape ]
	Since the bent solenoids introduce dispersion into the beam and the dipole fields compensate for this, the height and momentum of the muon beam at the stopping target can be controlled to a degree.
	Whilst in the optimisation of the muon beam dipoles of this thesis the target shape was held fixed, in principle some of the identified correlations might be different if the target design was allowed to vary simultaneously.  
	Such a study would clearly involve an enormous parameter space, so perhaps this is a study that could be best performed with a smarter machine-learning based approach.
\item[ Spectrometer's solenoidal field strength and aperture]
	Since the dispersive effect of a bent solenoid is inversely proportional to the solenoidal field strength, reducing the field in the spectrometer will likely improve the DIO--signal separation, reducing the hit rate.
	This could however also affect the signal acceptance since the trajectories will acquire larger envelopes.
	Increasing the spectrometer aperture size could compensate this, but then one faces an increase in the cost of the spectromter and the detector solenoid which would need to be considered.
\end{description}
